\renewcommand{\abstractname}{摘要}
\begin{abstract}

随着人类向地外空间拓展的步伐加速，月球殖民地的建设已从科幻愿景转变为可行的工程目标。月球殖民地管理机构（MCM）计划于2050年启动建设容纳\textbf{10万居民}的大规模月球栖息地，这一宏伟工程需要从地球向月球运输约\textbf{1亿公吨}建设材料——相当于建造超过1000座帝国大厦的总质量。如此空前的物流挑战亟需创新的运输解决方案。本文构建了一套\textbf{多模态星际物流优化框架}，系统评估太空电梯系统、传统火箭发射网络及其混合策略在成本、时间和环境影响三个维度的综合表现。

针对\textbf{太空电梯系统}，本文建立了三座银河港的运力模型。每座银河港配备两条10万公里石墨烯系绳连接地球赤道港口与顶点锚点，设计年运力\textbf{17.9万公吨}。考虑\textbf{2\%}年故障率和\textbf{5\%}系绳动力学运力折减后，系统有效年运力为\textbf{49.99万公吨}。电力驱动的升降机实现了\textbf{100美元/公斤}的超低运输成本和仅\textbf{0.1公吨CO$_2$/公吨载荷}的碳强度——较火箭运输减排\textbf{97.5\%}。纯电梯方案总成本\textbf{10万亿美元}，但建设周期长达\textbf{190年}。

针对\textbf{传统火箭发射网络}，本文整合了分布于美国（5处）、哈萨克斯坦、法属圭亚那、印度、中国和新西兰的10个发射基地，构建了全球化发射能力模型。基于2050年先进猎鹰重型火箭的技术预测，单次运载能力\textbf{100-150公吨}，年发射总量可达\textbf{3250次}。在\textbf{3\%}发射失败率约束下，有效年运力\textbf{39.41万公吨}，运输成本\textbf{1500美元/公斤}。纯火箭方案虽可直接地月运输，但总成本高达\textbf{150万亿美元}，建设周期\textbf{254年}，且产生\textbf{4亿公吨CO$_2$}排放。

本文核心贡献在于提出\textbf{运力比例分配优化模型}。该模型以最小化建设周期为目标，按两系统运力比例分配材料，确保并行运营的同步完成。求解得最优配置为：太空电梯承担\textbf{55.9\%}（5592万公吨），火箭承担\textbf{44.1\%}（4408万公吨）。混合方案实现合计年运力\textbf{89.4万公吨}，将建设周期压缩至\textbf{112年}（较纯电梯缩短\textbf{41\%}，较纯火箭缩短\textbf{56\%}），总成本\textbf{71.7万亿美元}，CO$_2$排放\textbf{1.82亿公吨}（较纯火箭减排\textbf{54.5\%}）。

敏感性分析验证了模型鲁棒性：电梯故障率从2\%升至12\%仅使工期延长\textbf{12.3\%}；火箭失败率同等变化仅影响\textbf{4.2\%}。殖民地运营期水需求分析显示，10万人年需\textbf{54.75万公吨}水资源，电梯运输成本\textbf{547.5亿美元/年}，仅为火箭运输的\textbf{6.7\%}——年节省超\textbf{7600亿美元}，凸显太空电梯对殖民地长期可持续运营的战略价值。本文建议MCM机构采用混合运输策略，优先投资太空电梯基础设施，为人类星际文明奠定物流基础。

\begin{keywords}
\textbf{太空电梯系统}；\textbf{月球殖民地物流}；\textbf{多模态运输优化}；\textbf{运力比例分配}；\textbf{星际基础设施}；\textbf{碳排放评估}
\end{keywords}

\end{abstract}
